\chapter*{前言}
\addcontentsline{toc}{chapter}{前言}

第一册把 C++ 程序追到二进制、ABI、汇编和可信测量；第二册把同一个程序继续放进现代 CPU、多核心、缓存、TLB、SIMD、线程、运行时、I/O 和分布式边界中。第三册从这里进入 AI 计算，但它不会把 AI 写成框架 API 清单。真正的问题不是“怎样调用一个模型”，而是：一个能推理开源 7B 级大模型的本地引擎，怎样把模型文件、tokenizer、权重张量、GEMM、attention、KV cache、量化参数、CPU/GPU kernel、调度器和测量报告组织成一条可解释的机器链路。

本册的贯穿项目是一台本地大模型推理引擎。目标不是玩具 MLP，而是能加载真实开源 7B 左右 decoder-only Transformer 权重，完成 prompt 到 token 流的推理，并且同一套 runtime 可以挂接 CPU 后端和 GPU 后端。书中仍会使用一个极小 MLP 作为第一块显微镜切片：它让矩阵乘、bias、activation、量化权重和 oracle 能被手工检查。但这个切片只服务主线，不取代主线。后续每一节都要回到真实引擎：模型格式怎样承载几十亿参数，prefill 和 decode 为什么有不同性能形态，KV cache 怎样决定显存/内存账本，量化怎样让 7B 模型装进本地机器，CPU 后端怎样一步步接近成熟框架，GPU 后端又怎样共享同一份算子合同。

\begin{keyidea}
第三册的主线是把大模型推理拆成可以验证的机器事实：张量是带形状、步长和生命周期的内存视图；算子是有输入合同和输出合同的 kernel；量化是整数表示、误差和吞吐之间的合同；KV cache 是跨 token 持续增长的状态；推理引擎是模型、内存、后端、调度和测量证据的组合。
\end{keyidea}

本册正文仍然以原理和工程证据为主。现有 \filepath{books/cpu-volume-3/labs/linux_cpu_inference/} 是第一阶段的最小切片；完整路线会继续扩展出模型加载、tokenizer、Transformer 层、KV cache、CPU kernel、GPU backend adapter、benchmark harness 和 correctness report。正文会使用伪代码、关键 C++ 片段、抽象汇编路径、数据布局表、流图和测量边界来解释机制；工程目录负责把这些机制落实到可构建、可测试、可 benchmark 的实现。
